% !TeX root = ../../../main.tex
\ChFourStarSection{Ethics statement} \label{ch04:app:ethics-statement}

\textsc{SurvHTE‐Bench} has significant positive potential for improving personalized medicine and clinical decision-making by enabling systematic evaluation of survival analysis methods under realistic assumption violations. By providing standardized benchmarks and practical guidance on when different estimators excel or fail, our work could accelerate the development of more reliable causal inference methods for high-stakes healthcare applications, ultimately supporting better patient outcomes through more informed treatment selection.

At the same time, our benchmark carries potential risks if misapplied. Practitioners may misinterpret benchmark results or place undue confidence in algorithmic decision-making, which could reduce necessary human oversight in clinical contexts. Moreover, although our study is methodological and does not involve human subjects directly, differences in estimator performance across demographic groups could exacerbate existing healthcare disparities if ignored. We therefore stress that our benchmark should not be used as a substitute for rigorous domain-specific validation, fairness assessment, or clinical trial evidence. 

All datasets used in this work are either publicly available synthetic or semi-synthetic datasets, or real-world datasets with proper access provisions (e.g., credentialed approval for MIMIC-IV). No personally identifiable information was used, and all data handling complies with the terms of use of the original sources. We encourage future applications of \textsc{SurvHTE‐Bench} to incorporate fairness audits, domain-specific validation, and appropriate safeguards to ensure responsible deployment.


